
\section{一些反例}
\label{sec:a1-5}

\begin{example}{A1.64}\label{ex:a1-64}
先给出一个与切丛的定义（定义~\ref{def:a1-14}）有关的例子。令
$X=[-2,2]$。取
\[
 V_1=\mathbb C\times[-1,2],\qquad
 E_1=\mathbb C,\qquad
 \Gamma_1=\{1\}.
\]
定义
\[
 s_1(z,t)=z^2,\qquad
 \psi_1(0,t)=t,\qquad
 \psi_1\colon s_1^{-1}(0)\longrightarrow X.
\]
再取
\[
 V_2=[-2,0],\qquad E_2=0,\qquad
 \Gamma_2=\{1\},\qquad s_2\equiv0,\qquad \psi_2(t)=t.
\]
坐标变换 $\phi_{12}$ 是
\[
 t\longmapsto(0,t),\qquad
 [-1,0]=V_{12}\longrightarrow
 \mathbb C\times[-1,0]\subset V_1.
\]
这些资料在 $X$ 上定义了一个仓西结构。法丛 $NV_{12}V_1$ 显然平凡，
因而同 $E_1/\widehat\phi_{12}(E_2)$ 同构。所以同构~\eqref{eq:a1-15}
存在，却不是由 $s$ 诱导的。

在 $\{2\}$ 处的重数为 $2$，而在 $\{-2\}$ 处的重数为 $1$；虚基本链
不同调于零。因此，由引理~\ref{lem:a1-31}，我们不能仅在边界之外修改
$X$ 的仓西结构，使修改后的结构具有切丛。
\end{example}

\begin{example}{A1.65}\label{ex:a1-65}
下面考察相容条件 (A1.6.2)。考虑轨形 $\mathbb C/\mathbb Z_2$，其中
$\mathbb Z_2$ 的作用为 $z\mapsto-z$。作为轨形，
$(\mathbb C/\mathbb Z_2)\setminus\{0\}$ 同
$\mathbb C\setminus\{0\}$ 同构。把它与
$(\mathbb C\cup\{\infty\})\setminus\{0\}$ 胶合，得到轨形 $X$。
这是一个\emph{坏轨形}的例子；换言之，它不是某个流形在离散群作用下
的\emph{全局商}。

$0$ 处的坐标由 $V_0/\Gamma_0=\mathbb C/\mathbb Z_2$ 给出。令
\[
 p_t=[\exp(2\pi t\sqrt{-1})]\in\mathbb C/\mathbb Z_2.
\]
以平凡的 $\Gamma$ 作用在一个小圆盘 $\operatorname{Int}D^2$ 上，得到
$p_t$ 的坐标图。从这个坐标图到 $V_0/\Gamma_0$ 的坐标变换为
\begin{equation}
 z\longmapsto \pm\exp(2\pi t\sqrt{-1})+z.
 \tag{A1.66}\label{eq:a1-66}
\end{equation}
上式中的正号或负号均可选取。

让 $t$ 从 $0$ 连续移动到 $1$。若不在相容条件 (A1.6.2) 中加入
$\gamma_{pqr}$，就不能改变式~\eqref{eq:a1-66} 中的正负号：一旦在某个
$t$ 处固定符号，之后便不允许改变。然而，$t=0$ 与 $t=1/2$ 对应同一个
点，这就产生了矛盾。也就是说，为了把这个例子纳入轨形的定义，必须在
定义~\ref{def:a1-5} 的条件 (A1.6.2) 中允许
$\gamma_{pqr}=-1$。
\end{example}

\begin{remark}{A1.67}\label{rem:a1-67}
轨形可以看作一个\emph{群胚}，即一个所有态射都是同构的范畴。（这是
经典事实，可以追溯到亚历山大·格罗滕迪克（Alexander Grothendieck）
开创叠理论的早期。正如第一、第四署名作者深谷贤治与小野薫在
\textbf{[FuOn99I]} 中所指出的，仓西结构理论是叠理论的一种类比。）

两个轨形之间的映射对应于这两个范畴之间的函子。在范畴论中，关键是要
区分两个对象、态射等彼此相等与彼此同构这两个概念。条件 (A1.6.2)
对应于两个函子等价（即同构）的条件，而取 $\gamma_{pqr}\equiv1$ 的同构
恰好对应于相等。

为了定义叠，需要胶合局部对象。胶合映射的相容条件，依定义要求胶合映射
的复合等价于第三个胶合映射；这正好同我们的定义中的条件 (A1.6.2)
一致。事实上，若使用\emph{2-范畴}\liftIndex{I375}的概念，上述叙述会更为系统。不过，
本书的应用不需要这种抽象，所以我们不用 2-范畴的语言，而把胶合条件
显式表述为等价类的相合条件。
\translatornote{源书把与复合 $\phi_{pq}\circ\phi_{qr}$ 相比较的
$\phi_{pr}$ 称为 ``the third composition map''。条件 (A1.6.2) 表明
$\phi_{pr}$ 是第三个直接胶合映射，而不是另一个复合；译文依其型别写作
“第三个胶合映射”。见疑误 A-0714。}
\end{remark}
